% Methods: governing equations and the curvilinear SBP-SAT discretization (as implemented in SeismicWavePC.jl)

\section{Governing equations}\label{sec:equations}

We solve the visco-elastic wave equation in velocity--stress form. Coordinates are $x$ (south to north in the 2D
sections), $y$, and $z$ (elevation, positive up). In 2D (P-SV, plane strain in the $x$--$z$ plane) the unknowns are the
particle velocities $v_x, v_z$ and the stresses $\sigma_{xx}, \sigma_{zz}, \sigma_{xz}$; in 3D also $v_y$, $\sigma_{yy}$,
$\sigma_{xy}$, $\sigma_{yz}$. Conservation of momentum reads
\begin{equation}
  \rho\,\frac{\partial v_i}{\partial t} = \frac{\partial \sigma_{ij}}{\partial x_j},
  \label{eq:momentum}
\end{equation}
with the density $\rho$ (summation over repeated indices). No body forces are applied in this study: the waves enter
through the boundary of an injection box (Sec.~\ref{sec:injection}).

\paragraph{Attenuation.} Constant $Q$ is modelled with a generalized Zener body of $L = 3$ relaxation mechanisms
($\tau$-method of Blanch et al., 1995, as in OpenSWPC; Maeda et al., 2017). With the strain rates
$\dot\varepsilon_{ij} = \tfrac12(\partial_j v_i + \partial_i v_j)$, the P-wave and shear moduli $M = \lambda + 2\mu$ and
$\mu$, and the memory variables $r^{(l)}_{ij}$, $l = 1,\dots,L$,
\begin{align}
  \frac{\partial \sigma_{ij}}{\partial t} &=
     \big[M(1+\tau_p) - 2\mu(1+\tau_s)\big]\,\dot\varepsilon_{kk}\,\delta_{ij}
     + 2\mu(1+\tau_s)\,\dot\varepsilon_{ij} + \sum_{l=1}^{L} r^{(l)}_{ij},
  \label{eq:stress}\\
  \frac{\partial r^{(l)}_{ij}}{\partial t} &= -\frac{1}{\tau_{\sigma l}}\Big( r^{(l)}_{ij}
     + \frac{1}{L}\Big[\big(M\tau_p - 2\mu\tau_s\big)\dot\varepsilon_{kk}\,\delta_{ij} + 2\mu\tau_s\,\dot\varepsilon_{ij}\Big]\Big).
  \label{eq:memory}
\end{align}
The relaxation times $\tau_{\sigma l}$ are spaced logarithmically between $f_{\min}$ and $f_{\max}$ of the simulated band,
$\tau_p = L\zeta/Q_p$ and $\tau_s = L\zeta/Q_s$, where $\zeta$ makes $Q$ constant over the band (Blanch et al., 1995). The
moduli are scaled so that the phase velocities at the reference frequency equal the given $v_p$ and $v_s$. Without
attenuation ($Q = \infty$), $\tau_p = \tau_s = 0$ and Eq.~\eqref{eq:stress} is Hooke's law. The memory variables are
integrated with the Crank--Nicolson rule.

\section{Discretization: curvilinear SBP-SAT finite differences}\label{sec:discretization}

\paragraph{Terrain-following coordinates.} The physical domain is mapped onto a rectangle with coordinates $(\xi,\eta)$
(3D: $(\xi,\psi,\eta)$). The grid columns are vertical and the top grid line follows the topography:
\begin{equation}
  x = \xi, \qquad (y = \psi), \qquad z = Z(\xi,\eta), \qquad
  J = \frac{\partial Z}{\partial \eta} > 0, \quad J_x = -\frac{\partial Z}{\partial \xi}, \quad (J_y = -\frac{\partial Z}{\partial \psi}).
\end{equation}
The chain rule gives the derivatives
\begin{equation}
  \frac{\partial}{\partial x} = \frac{\partial}{\partial \xi} + \frac{J_x}{J}\frac{\partial}{\partial \eta}, \qquad
  \frac{\partial}{\partial z} = \frac{1}{J}\frac{\partial}{\partial \eta},
  \label{eq:chain}
\end{equation}
and, with the metric identity $\partial_\xi J + \partial_\eta J_x = 0$, the momentum equation in conservative form,
\begin{equation}
  \rho J\,\frac{\partial v_i}{\partial t} = \frac{\partial}{\partial \xi}\big(J\sigma_{ix}\big)
     + \frac{\partial}{\partial \eta}\big(J_x\sigma_{ix} + \sigma_{iz}\big)
  \qquad (i = x, z).
  \label{eq:momentum_curv}
\end{equation}
The flux $T_i = J_x\sigma_{ix} + \sigma_{iz}$ through a grid line $\eta = $ const is the traction on that line times
$|\nabla\eta|\,J$; at the free surface $T_i = 0$ is the traction-free condition.

\paragraph{SBP operators.} Along each coordinate the derivative is approximated with \emph{dual-pairing upwind}
summation-by-parts operators $D_+$, $D_-$ (Mattsson, 2017) with a diagonal, positive norm (quadrature) $H$. They satisfy
\begin{equation}
  H D_+ + \big(H D_-\big)^{\mathsf T} = B = e_N e_N^{\mathsf T} - e_1 e_1^{\mathsf T},
  \label{eq:sbp}
\end{equation}
the discrete analogue of integration by parts, $\int u\,w' + \int u'\,w = [u\,w]$. We use the order-7 operators with the
boundary-optimized closures of Mattsson, Niemelä \& Winters (2027), from SummationByPartsOperators.jl (Ranocha, 2021). The
divergence in the momentum equation uses $D_-$, the gradient in the strain rates $D_+$:
\begin{align}
  \rho J\,\frac{d v_i}{dt} &= D_-^{\xi}\big(J\sigma_{ix}\big) + D_-^{\eta}\big(J_x\sigma_{ix} + \sigma_{iz}\big)
     + \mathrm{SAT}_i , \label{eq:semi_momentum}\\
  \dot\varepsilon_{xx} &= D_+^{\xi} v_x + \frac{J_x}{J} D_+^{\eta} v_x, \qquad
  \dot\varepsilon_{zz} = \frac{1}{J} D_+^{\eta} v_z, \qquad
  2\dot\varepsilon_{xz} = \frac{1}{J} D_+^{\eta} v_x + D_+^{\xi} v_z + \frac{J_x}{J} D_+^{\eta} v_z ,
  \label{eq:semi_strain}
\end{align}
and the stresses follow from Eqs.~\eqref{eq:stress}--\eqref{eq:memory} node by node. All fields are collocated at the
grid nodes.

\paragraph{Boundary conditions (SAT).} Boundary and interface conditions are imposed weakly with simultaneous
approximation terms (Carpenter et al., 1994). At the free surface (the top row $\eta_N$) the SAT removes the flux,
\begin{equation}
  \mathrm{SAT}_i = -\,H_\eta^{-1} e_N\, e_N^{\mathsf T}\big(J_x\sigma_{ix} + \sigma_{iz}\big),
  \label{eq:sat_free}
\end{equation}
so that the discrete traction vanishes in the sense of the SBP norm. At the sides and the bottom the traction is replaced by
the impedance condition $\sigma_{ij} n_j = -Z v_i$ (with $Z = \rho v_p$ for the normal and $\rho v_s$ for the tangential
component), integrated implicitly with the trapezoidal rule, so that it does not restrict the time step; absorbing layers
(below) take out the remaining waves. Material interfaces, where present, are grid lines between blocks; there half the
traction mismatch enters the momentum equations and half the velocity mismatch the strain rates of both sides.

\paragraph{Energy estimate.} With Eq.~\eqref{eq:sbp}, the boundary terms of $D_-^{\eta}$ in Eq.~\eqref{eq:semi_momentum}
cancel against Eq.~\eqref{eq:sat_free}. For the elastic case the semi-discrete energy
\begin{equation}
  E = \frac12 \sum_{\text{nodes}} H_\xi H_\eta \Big( \rho J\, |v|^2 + J\,\sigma : \mathsf{C}^{-1} \sigma \Big)
\end{equation}
is therefore conserved by the interior scheme and the free surface, and it decreases at the impedance boundaries: the
scheme is stable for any grid and any (smooth) topography. The absorbing layers below are not covered by this estimate. The same identity, $H A = -(H J G)^{\mathsf T}$ between the
discrete divergence $A$ and gradient $G$ including the SAT terms, makes the discrete adjoint exact (used for misfit
gradients; not needed here).

\paragraph{Time stepping.} The velocities live at half time steps (leapfrog):
\begin{equation}
  v^{n+1/2} = v^{n-1/2} + \Delta t\, \mathcal{A}\big(\sigma^{n}\big), \qquad
  \sigma^{n+1} = \sigma^{n} + \Delta t\, \mathcal{G}\big(v^{n+1/2}\big),
\end{equation}
with $\mathcal{A}$ and $\mathcal{G}$ the right-hand sides of Eqs.~\eqref{eq:semi_momentum} and
\eqref{eq:stress}--\eqref{eq:semi_strain}. The scheme is stable for $\Delta t < 2/\sqrt{\lambda_{\max}}$, with
$\lambda_{\max}$ the largest eigenvalue of $\mathcal{A}\mathcal{G}$, found by power iteration; we use 90\,\% of this limit.

\paragraph{Absorbing layers.} The domain is surrounded by complex-frequency-shifted perfectly matched layers (CFS-PML)
implemented with auxiliary differential equations (Zhang \& Shen, 2010), applied to the $\xi$ and $\eta$ parts of the SBP
operators.

\section{Plane-wave injection}\label{sec:injection}

The incoming plane waves are injected with the total-field/scattered-field method: inside a box (sides and bottom; it
reaches up to the free surface) the solver computes the total field, outside only the field scattered by the structure in
the box, which the absorbing layers remove. The incoming field of the horizontally layered background, i.e.\ the plane
wave together with its free-surface reflection and the reverberations in the layers, is computed exactly in the frequency
domain (global matrix of up- and down-going waves). The velocity gradient is represented by one layer per mesh row.
Every time step, the finite-difference stencils that cross the box boundary are corrected with this field. The angle of
incidence is defined in the half-space below the model:
\begin{equation}
  p = \frac{\sin\theta_0}{v_0}, \qquad \sin\theta(z) = p\, v(z),
  \label{eq:snell}
\end{equation}
so that the waves bend towards the vertical as the velocity decreases upwards. To reproduce an incidence of $45^\circ$ at
$z = 0$ (the paper's tilted source plane), $\theta_0 = \arcsin\!\big(\sin 45^\circ\, v_p(-2\,\mathrm{km})/v_p(0)\big)
\approx 56.8^\circ$. One wave type (P, or S polarized along $x$ or $y$) is injected per run, and the runs are summed.

\section{Stress measure}\label{sec:dq}

Following the paper's supplement, the dynamic stress is measured by the deviatoric stress
\begin{equation}
  \mathrm{d}Q = \sqrt{J_2}, \qquad J_2 = \tfrac12\, s_{ij} s_{ij}, \qquad
  s_{ij} = \sigma_{ij} - P\,\delta_{ij}, \qquad P = \tfrac13\big(\sigma_{xx} + \sigma_{yy} + \sigma_{zz}\big),
  \label{eq:dq}
\end{equation}
where in 2D the out-of-plane stress of plane strain, $\sigma_{yy} = \lambda(\varepsilon_{xx} + \varepsilon_{zz})
= \nu(\sigma_{xx} + \sigma_{zz})$, is included. Figure~7 shows, at every node, the maximum over time of $\mathrm{d}Q$.
